% \iffalse meta-comment
%
% hit-options.dtx 是各个类共用的选项声明与判断助手
%
% \fi
%
% \section{共享选项层}
%
% \changes{v3.2a}{0000/00/00}{声明选项的四个助手与一族判断函数抽出来给各个类共用，
%   原先每个类各写一遍，十五个定义逐字节相同}
% 声明类选项的助手，以及排版代码里用来读标志位的一族判断函数。原先每个类的
% \texttt{*-options} 模块里各写一遍，十五个定义一字不差。
%
% 这个文件只放定义，不执行任何依赖类状态的东西，所以排在类自己的模块之前没关系。
% 各个类照旧在自己的 \texttt{*-options} 里调用 \cs{hit@boolopt} 这些声明选项。
%
%    \begin{macrocode}
%<*shared-options>
%    \end{macrocode}
%
% 声明选项
%
% \cs{hit@boolopt}\marg{名}\marg{默认} 声明一个布尔选项，标志位是
% \cs{g\_\_hit\_}\meta{名}\cs{\_bool}；\cs{hit@stringopt}\marg{名} 声明一个存字符串的选项，
% 存进 \cs{hit@}\meta{名}。键名带连字符时不能直接拿来拼控制序列，
% \cs{hit@stringopt@as}\marg{键名}\marg{宏名} 把两者分开。
%
% \cs{hit@choiceopt}\marg{键名}\marg{前缀}\marg{取值表}\marg{默认} 声明一组互斥取值，
% 每个取值一个标志位，选中的置真、其余置假，取值不在表里就报错。
%    \begin{macrocode}
\msg_new:nnn { hithesis } { bad-option-value }
  { 选项~#1~不接受取值~#2。 }
\cs_new_protected:Npn \hit@boolopt #1#2
  {
    \bool_new:c { g__hit_#1_bool }
    \str_if_eq:nnT {#2} { true } { \bool_gset_true:c { g__hit_#1_bool } }
    \keys_define:nn { hit / class }
      { #1 .bool_gset:c = { g__hit_#1_bool } , #1 .default:n = { true } }
  }
\cs_new_protected:Npn \hit@stringopt #1 { \hit@stringopt@as {#1} {#1} }
\cs_new_protected:Npn \hit@stringopt@as #1#2
  {
    \tl_new:c { hit@#2 }
    \keys_define:nn { hit / class } { #1 .tl_gset:c = { hit@#2 } }
  }
\cs_new_protected:Npn \hit@choiceopt #1#2#3#4
  {
    \clist_map_inline:nn {#3} { \bool_new:c { g__hit_#2##1_bool } }
    \tl_if_empty:nF {#4} { \bool_gset_true:c { g__hit_#2#4_bool } }
    \keys_define:nn { hit / class }
      {
        #1 .code:n =
          {
            \clist_map_inline:nn {#3}
              { \bool_gset_false:c { g__hit_#2####1_bool } }
            \bool_if_exist:cTF { g__hit_#2##1_bool }
              { \bool_gset_true:c { g__hit_#2##1_bool } }
              { \msg_error:nnnn { hithesis } { bad-option-value } {#1} {##1} }
          } ,
      }
  }
%    \end{macrocode}
%
% 读标志位的判断函数
%
% 排版代码里大量出现「按某个标志位二选一」，写成 \cs{bool\_if:cTF} 要带上
% \cs{g\_\_hit\_}\ldots\cs{\_bool} 这一长串，所以各起了名字。参数是在 \pkg{expl3}
% 之外记号化的，里面的字面空格照原样保留——这正是它们存在的另一个理由，
% 含排版内容的判断不必整段搬进 \pkg{expl3}。
%    \begin{macrocode}
\cs_new:Npn \hit@if #1#2#3 { \bool_if:cTF { g__hit_#1_bool } {#2} {#3} }
\cs_new:Npn \hit@ifT #1#2 { \bool_if:cT { g__hit_#1_bool } {#2} }
\cs_new:Npn \hit@ifand #1#2#3#4
  {
    \bool_lazy_and:nnTF
      { \bool_if_p:c { g__hit_#1_bool } } { \bool_if_p:c { g__hit_#2_bool } } {#3} {#4}
  }
\cs_new:Npn \hit@iffontset #1#2#3
  { \str_if_eq:eeTF { \hit@fontset } {#1} {#2} {#3} }
\cs_new_protected:Npn \hit@ifnotempty #1#2 { \tl_if_empty:NF #1 {#2} }
\cs_new:Npn \hit@ifwider #1#2#3#4 { \dim_compare:nNnTF {#1} > {#2} {#3} {#4} }
\cs_new:Npn \hit@ifwiderdim #1#2#3#4 { \dim_compare:nNnTF {#1} > {#2} {#3} {#4} }
\cs_new:Npn \hit@ifmore #1#2#3 { \int_compare:nNnT {#1} > {#2} {#3} }
\cs_new_protected:Npn \hit@ifsame #1#2#3#4
  {
    \group_begin:
      \protected@edef \l__hit_cmpa_tl {#1}
      \protected@edef \l__hit_cmpb_tl {#2}
      \tl_if_eq:NNTF \l__hit_cmpa_tl \l__hit_cmpb_tl
        { \group_end: #3 }
        { \group_end: #4 }
  }
%    \end{macrocode}
%
% style/glue 的伸缩量
%
% \option{style/glue} 开时给标题前后距与行距一个伸缩量，关时是固定值。
% 这两个函数把原地重复过二十多次的判断收成一处。
%    \begin{macrocode}
\cs_new:Npn \hit@glueskip #1#2
  { \bool_if:NTF \g__hit_glue_bool { #1~plus~#2~minus~0bp } { #1 } }
\cs_new:Npn \hit@gluesize #1#2#3
  {
    \bool_if:NTF \g__hit_glue_bool
      { \fontsize {#1} { #2~plus~#3 } } { \fontsize {#1} {#2} }
  }
%    \end{macrocode}
%
%    \begin{macrocode}
%</shared-options>
%    \end{macrocode}
%
% \Finale
